% ch01 -- The Bit as a Physical Object

\epigraph{A bit is not a thing. It is a decision.}{}

\section{What is a Bit Really?}

The intuitive answer: a bit is a $0$ or a $1$. That is a
descriptive decision, not a physical statement.

The physical statement is different: A bit is a physical
system with two distinguishable macro-regions in phase space. \status{B}
Whether one calls these regions $\{0, 1\}$, $\{\uparrow, \downarrow\}$ or
$\{\text{left, right}\}$ changes nothing about the physics. The naming
is a descriptive decision of the observer --- not a law of nature.

\begin{laierspur}
Imagine two glasses on a table. One is full, one is empty.
These are two distinguishable physical states --- a bit,
before you call it that. Whether you say "full = 1, empty = 0" or vice versa,
does not change the physics. The naming is your decision.
\end{laierspur}

\begin{keyresult}[Core Distinction]
\textbf{Physical Region} $\neq$ \textbf{Description of the Region} \status{SETZUNG}

The region in phase space exists independently of how one divides it.
The naming as "bit" is an external interpretation.
\end{keyresult}

\section{Multiple Realization: The Same Bit, Different Carriers}

A bit with the value $1$ can be realized by fundamentally different physical
systems: a CMOS transistor at 5\,V, a magnetic
spin $\uparrow$, a circularly polarized photon, an abacus bead on the left of the
working field, a nanomagnet pillar with orientation $+z$.

The information content is identical in all cases: 1 bit. The energies
vary by many orders of magnitude. It follows: Energy is a
property of the \textbf{carrier}, not of the information. \status{B}

\begin{laierspur}
A letter contains the same information whether you send it by airmail or
by e-mail. The energy of transmission is completely different ---
the information content is identical. Energy clings to the carrier,
not to the information.
\end{laierspur}

\section{The Three Levels of Description}

In the FFGFT, three levels are strictly separated (Doc.\,A272): \status{B}

\paragraph{Level 1 --- Construction.}
The engineer determines during chip design which region operations
physically exist --- bijective or non-bijective. This is cast in silicon,
unchangeable during operation. This decision determines
the thermodynamic cost structure --- not the program.

\paragraph{Level 2 --- State Transition.}
During operation, state transitions take place in a clocked manner.
Landauer costs arise \emph{only} for non-bijective region mappings.
The clock frequency, not the computational operation, determines the entropy flow.

\paragraph{Level 3 --- Description.}
The program chooses which hardware is used. That is
resource management. The thermodynamic costs still arise ---
determined by construction and switching operation, independent of
what the description level does with the results.

\section{How Many States Are in a Region?}

A colloidal particle in a 1-\textmu m trap contains physically about
$7{,}9 \cdot 10^9$ orthogonal quantum states. If one names roughly two zones
"left" and "right," one describes this region with 1 logical bit.
If one divides it into 1000 zones, one describes the same region with ${\sim}10$
logical bits. The physical region is identical in both cases ---
the depth of description is a choice of the observer. \status{B}

The counting basis is always $h^f$: the fundamental discretization of the
$f$-dimensional phase space, independent of naming and depth of description.

\section{Bit, p-Bit, Qubit: A Scale}

The classical bit operates in a deep potential well $U \gg k_B T$ ---
macroscopically stable, for years. The qubit requires metastable coherence
in the millisecond range at mK temperatures. In between lies the
\textbf{p-bit}: a magnetic nanostructure in which the torsion barrier
becomes thermally surmountable.

The FFGFT shows that the natural scale for p-bits is geometrically
predetermined: Level $N = 8$ of the $\xi$ hierarchy yields
$L_8 \approx 21{,}6\,\text{nm}$ --- exactly the transistor limit size. \status{K Doc.\,184}

\begin{laierspur}
Three operating points, one scale. The classical bit is like a heavy
ball deep in a valley --- stable for years. The qubit is a ball
on a mountain peak, held there by deep cooling.
The p-bit lies in between: the valley is flat enough that heat occasionally
kicks the ball to the other side. The FFGFT geometry predicts
where this boundary lies: at 22 nanometers.
\end{laierspur}

These three regimes --- classical bit, p-bit, qubit --- are not
separate worlds. They are three operating points on the same physical
scale. \status{B}